\documentclass[10pt,a4paper]{article}
\usepackage{amsmath,amssymb,amsthm,graphicx}
\newcommand{\rk}{\operatorname{rk}}
\newcommand{\xc}{\operatorname{xc}}
\newcommand{\sep}{\operatorname{sep}}
\newcommand{\SGN}{\operatorname{SGN}}
\newcommand{\sgn}{\operatorname{sgn}}
\newcommand{\conv}{\operatorname{conv}}
\newcommand{\RR}{\mathbb R}
\newenvironment{originalresult}[1]{\par\medskip\noindent\textbf{#1.}\itshape}{\par\medskip}
\title{Hard submatrices for Boolean 1-partition number}
\author{Pavel Hrubeš}
\date{doi:10.4230/LIPIcs.CCC.2024.13; CC BY 4.0}
\begin{document}
\maketitle
\noindent\textit{Editorial note.} Faithful transcription from official CCC 2024 CC BY 4.0 PDF; original notation, indexing and dimensions retained. Author TeX was not located. No author proof is replaced or extended; standard cited prior results are retained as prerequisites, and bibliography entries keep their original numbers.\par\medskip
\noindent\textit{Allowed standard prerequisites (editorial).}\begin{itemize}
\item Elementary rank and binomial-coefficient bounds. Author source Sections 3-4
\end{itemize}\medskip
\section*{3 A combinatorial lemma (author source pp. 13:3--13:4)}
Both Theorems 1 and 2 rely on a simple combinatorial lemma.
\begin{originalresult}{Lemma 3}
Let $\mathcal A\subseteq 2^{[n]}$ be a family of subsets of $[n]$. Assume that $1\le k\le n$ is such that every $k$-element subset of $[n]$ is contained in some $A\in\mathcal A$. Then there exists a subfamily $\mathcal A'\subseteq\mathcal A$ of size $|\mathcal A'|\le O(|\mathcal A|^{1/k}\log(n/k))$ with $\bigcup\mathcal A'=[n]$. In particular, if $|\mathcal A|\le 2^k$ then $|\mathcal A'|\le O(\log n)$.
\end{originalresult}
\begin{proof}
Assume that $|\mathcal A|\le a^k$. Let $t$ be the size of a largest set in $\mathcal A$. Then we have
\[\binom nk\le a^k\binom tk.\]
Hence $t\ge\frac n{ea}$, using the estimates $(\frac nk)^k\le\binom nk$, $\binom tk\le(\frac{et}k)^k$. Take some $A_0\in\mathcal A$ of size $t$. Let
\[\mathcal A_1:=\{A\setminus A_0:A\in\mathcal A\}.\]
Then every subset of $U_1:=[n]\setminus A_0$ of size at most $k$ is contained in some member of $\mathcal A_1$. The size of $U_1$ is at most $n(1-\frac1{ea})$. Similarly, take a largest set $A_1$ from $\mathcal A_1$ and obtain a new family $\mathcal A_2\subseteq 2^{U_2}$ on $U_2:=U_1\setminus A_1$. After $s$ steps, the size of $U_s$ is at most $n(1-\frac1{ea})^s$ and after $s\le O(a\log(n/k))$ steps we have $|U_s|\le k$. This guarantees that the largest set in $\mathcal A_s$ is $U_s$ itself and $[n]=\bigcup_{i=0}^s A_s$. By construction, every $A_i$ is contained in some element of the original family $\mathcal A$.
\end{proof}

\section*{Definitions and notation (author source p. 13:3)}
Given a Boolean matrix $M\in\{0,1\}^{m\times n}$, let us define its 1-partition number, $\chi_1(M)$, as the smallest $s$ such that $M$ can be written as
\[M=LR,\quad\text{with }L\in\{0,1\}^{m\times s},\ R\in\{0,1\}^{s\times m},\]
where the operations are over $\RR$. The definition emphasizes the analogy with $\rk_+$, and $\chi_1$ is also sometimes referred to as binary rank. On the other hand, the phrase ``partition number'' comes from communication complexity. The name is justified: it is easy to see that $\chi_1(M)$ equals the smallest $s$ such that the 1-entries of $M$ can be partitioned into $s$ 1-monochromatic rectangles (i.e., rank-one Boolean matrices). Finally, when $M$ is viewed as the adjacency matrix of a bipartite graph, $\chi_1(M)$ also appears under the name biclique partition number [3].
\paragraph{Notation.} All logarithms are in base 2 and $[n]:=\{1,\ldots,n\}$.
\section*{4 1-Partition number (author source pp. 13:4--13:5)}
In this section, we prove Theorem 2.
Let $M$ be an $m\times n$ matrix with rows indexed by $[n]=\{1,\ldots,n\}$. Given $A\subseteq[n]$, $M_A$ denotes the submatrix obtained by removing the rows outside of $A$ from $M$. Observe that\footnote{If $A_1,A_2$ are disjoint, this is quite obvious. Otherwise consider $A_1,A_2\setminus A_1$.}
\[\chi_1(M_{A_1\cup A_2})\le\chi_1(M_{A_1})+\chi_1(M_{A_2}),\tag{1}\]
and so $\chi_1(M_A)$ can be viewed as a subadditive measure on $[n]$ whenever $M$ is fixed.
If a matrix $M$ has rank $r$, its rows are a linear combination of a subset of $r$ rows of $M$. This means that every column of $M$ is determined by a fixed subset of $r$ coordinates. If $M$ is Boolean, this leads to the following useful fact:
if $M$ has distinct columns then $n\le 2^{\rk(M)}$ (similarly for rows).
\begin{originalresult}{Lemma 5}
Let $M$ be an $m\times n$ Boolean matrix of rank $r$. Given $s\in[n]$, let $\mathcal A$ be the collection of maximal subsets $A\subseteq[n]$ with $\chi_1(M_A)\le s$ (i.e., $\chi_1(M_A)\le s$ and $\chi_1(M_{A'})>s$ for every $A'\supsetneq A$). Then $|\mathcal A|\le 2^{(r+s)^2}$.
\end{originalresult}
\begin{proof}
Let $v_1,\ldots,v_n\in\RR^m$ be the columns of $M$. Given $L\in\{0,1\}^{m\times s}$, let
\[L^*:=\{i\in[n]:\exists y\in\{0,1\}^s\ v_i=Ly\}.\]
Let $\mathcal L:=\{L^*:L\in\{0,1\}^{m\times s}\}$.
We claim that $\mathcal A\subseteq\mathcal L$. If $\chi_1(M_A)\le s$, we can write $M_A=LR$ with $L\in\{0,1\}^{m\times s}$ and $R\in\{0,1\}^{s\times|A|}$. This means that every $v_i$, $i\in A$, is a Boolean linear combination of the columns of $L$ and $A\subseteq L^*$. Furthermore, if $A$ is maximal, we must have $A=L^*$.
We now want to estimate the size of $\mathcal L$. The set $L^*$ consists of indices $i\in[n]$ so that there exists $x\in\RR^n,y\in\RR^s$ satisfying
\[Mx-Ly=0\tag{2}\]
such that $y\in\{0,1\}^s$ and $x$ is the $i$-th unit vector. Since $M$ has rank $r$ and $L$ has rank at most $s$, the system (2) is equivalent to a subsystem of $t:=\min((s+r),m)$ equations. These correspond to rows of the matrix $(M,L)$. Hence, in order to determine $L^*$, it is sufficient to specify a $t$-element subset of $[m]$ together with the $t\times s$ submatrix of $L$. This gives the estimate
\[|\mathcal L|\le\binom mt2^{ts}\le2^{t(s+\log m)}.\]
Finally, we can assume that $M$ has distinct rows and so $\log m\le r$, obtaining the bound $2^{(r+s)^2}$.
\end{proof}

\begin{originalresult}{Theorem 2 (restated; author source p. 13:5)}
Let $M$ be an $m\times n$ Boolean matrix with $n\ge2$. Then for every $k\le n$, $M$ contains an $m\times k$ submatrix of $k$ columns with 1-partition number $\Omega(\min(\sqrt{k},\frac{\chi_1(M)}{\log n}))$.
\end{originalresult}

\begin{proof}
Let $r$ be the rank of $M$. We will assume $r\le\frac{k^{1/2}}2$. Otherwise, observe that $M$ contains a full rank $r\times r$ submatrix, $\chi_1$ is lower-bounded by rank, and the conclusion of the theorem follows.
Let $s$ be the maximum $\chi_1(M_A)$ over all $A\subseteq[n]$ of size $k$. Let $\mathcal A$ be the family from the previous lemma. If $|\mathcal A|\ge2^k$, we have $2^k\le2^{(s+r)^2}$ and therefore $s\ge\frac{k^{1/2}}2$ from the assumption on $r$.
Assume $|\mathcal A|\le2^k$. By Lemma 3, there exists a subfamily $\mathcal A'\subseteq\mathcal A$ of size $O(\log n)$ which covers $[n]$. Using (1), this implies $\chi_1(M)\le O(s\log n)$ and so $s\ge\Omega(\chi_1(M)/\log n)$.
\end{proof}

\begin{thebibliography}{99}

\bibitem[1]{source1} S. Basu, R. Pollack, and M.F. Roy. Algorithms in real algebraic geometry. Springer-Verlag, 2006.
\bibitem[2]{source2} G. Braun, S. Fiorini, S. Pokutta, and D. Steuer. Approximation limits of linear programs (beyond hierarchies). Math. of Operations Research, 40(3), 2015.
\bibitem[3]{source3} S. Chandran, D. Issac, and A. Karrenbauer. On the parameterized complexity of biclique cover and partition. In International Symposium on Parameterized and Exact Computation, 2016.
\bibitem[7]{source7} P. Hrube\v{s}. On $\epsilon$-sensitive monotone computations. Computational Complexity, 29(2), 2020.
\bibitem[8]{source8} P. Hrube\v{s}. On the complexity of computing a random Boolean function over the reals. Theory of Computing, 16(9):1--12, 2020.
\bibitem[9]{source9} P. Hrube\v{s} and N. Talebanfard. On the extension complexity of polytopes separating subsets of the Boolean cube. Disc. and Comp. Geom., 70(1), 2022.
\bibitem[18]{source18} Thomas Rothvo\ss. Some 0/1 polytopes need exponential size extended formulations. CoRR, abs/1105.0036, 2011. arXiv:1105.0036.
\bibitem[20]{source20} Mihalis Yannakakis. Expressing combinatorial optimization problems by linear programs. Journal of Computer and System Sciences, 43(3):441--466, 1991.

\end{thebibliography}
\end{document}
